%% Elements of Nested Fibrational Cosmology
%% Book V: Descents, Projections, and Certified Branch Legitimacy
%% Final-Pass Canonical Draft
%%
%% Status legend:
%%   [U]  Unconditional  — proved from standing definitions and prior Unc results
%%   [C]  Conditional    — depends on explicitly declared hypotheses
%%   [D]  Definition-structural — structural, no proof obligation
%%   [R]  Remark only    — expository, non-load-bearing
%%   [O]  Open obligation — not yet proved

\documentclass[12pt,oneside]{amsbook}

\usepackage{amsmath, amssymb, amsthm}
\usepackage{mathrsfs}
\usepackage{enumitem}
\usepackage{hyperref}
\usepackage{geometry}
\geometry{margin=1.2in}

%% ---------------------------------------------------------------
%% Theorem environments
%% ---------------------------------------------------------------

\theoremstyle{plain}
\newtheorem{theorem}{Theorem}[section]
\newtheorem{lemma}[theorem]{Lemma}
\newtheorem{proposition}[theorem]{Proposition}
\newtheorem{corollary}[theorem]{Corollary}

\theoremstyle{definition}
\newtheorem{definition}[theorem]{Definition}
\newtheorem{standingrule}[theorem]{Standing Rule}

\theoremstyle{remark}
\newtheorem{remark}[theorem]{Remark}
\newtheorem{scholium}[theorem]{Scholium}

%% Status tag macro
\newcommand{\status}[1]{\textup{[#1]}}

%% ---------------------------------------------------------------
%% Notation
%% ---------------------------------------------------------------
%%  From Books I--IV (all freely used):
%%    \Omega, \Sigma_\infty, \delta, \Delta, O(U), C_{\rho^*}(U)
%%    \mathrm{Reg}, \mathrm{POT}, F_R, \tau_\iota, \mathbf{U}, K^*, T^*, \pi
%%
%%  New in Book V:
%%    \mathcal{C}           target category of a branch
%%    F: \mathrm{POT} \to \mathcal{C}   realization functor
%%    p: \mathbf{U} \to E   branch projection
%%    E                     endpoint datum
%%    W                     legitimacy witness
%%    \mathrm{Obs}^{\mathrm{stab}}  algebra of stabilized admissible observables

\begin{document}

\title{Book V\\[6pt]Descents, Projections, and Certified Branch Legitimacy}
\author{Elements of Nested Fibrational Cosmology}
\date{}
\maketitle

\tableofcontents

\bigskip

%% ---------------------------------------------------------------
\section*{V.0\quad Purpose}
%% ---------------------------------------------------------------

Book~V is the descent and governance book of the canonical spine.

Book~IV proved that the universal source object $\mathbf{U}$ exists,
is bi-universal, and carries the canonical layered structure
$T^* \hookrightarrow K^* \hookrightarrow \mathbf{U}$.  The program
now faces a new question: given that there is one universal source,
how can it lawfully determine multiple downstream branches without
collapsing them into arbitrary reinterpretations or permitting hidden
imports?

Book~V answers this question by defining and characterizing
\emph{lawful descent}.  Its central theorem --- the Universal Branch
Legitimacy Theorem --- gives a necessary and sufficient condition for
a proposed downstream branch to count as a lawful descendant of
$\mathbf{U}$.  The five conditions are: explicit branch candidate
intake, certified image of $\mathbf{U}$ as the endpoint datum,
endpoint visibility, admission of a legitimacy witness, and exclusion
of hidden-channel supplementation.

Book~V is explicitly a \emph{law-of-branchhood} book.  It does not
close any specific endpoint program.  It does not prove Yang--Mills,
Navier--Stokes, gravity, or any other specific result.  Its purpose is
to make the structure of lawful descent precise so that branch books
cannot claim more than their declared descent machinery licenses.

%% ---------------------------------------------------------------
\section{Imports and Standing Rules}
%% ---------------------------------------------------------------

\textbf{Imports.}  All certified results of Books~I--IV.  In
particular: the bi-universal property of $\mathbf{U}$
(Book~IV, Corollary~IV.5.1), the canonical layered structure
$T^* \hookrightarrow K^* \hookrightarrow \mathbf{U}$
(Book~IV, Corollary~IV.6.1), and the result that every lawful
downstream realization factors through $\mathbf{U}$
(Book~IV, Corollary~IV.5.3).

\begin{standingrule}[\status{D}]\label{sr:source-descent-only}
A branch is legitimate only if its endpoint content is sourced through
$\mathbf{U}$ and declared descent machinery.
\end{standingrule}

\begin{standingrule}[\status{D}]\label{sr:endpoint-visibility}
No endpoint feature may be counted as branch structure unless it is
branch-visible in certified terms.
\end{standingrule}

\begin{standingrule}[\status{D}]\label{sr:no-hidden-supplementation}
A downstream claim may not import undeclared structure under the guise
of interpretation, intuition, or endpoint convenience.
\end{standingrule}

\begin{standingrule}[\status{D}]\label{sr:legitimacy-before-closure}
A branch must first be shown lawful before any branch-specific closure
theorem has force.
\end{standingrule}

\begin{standingrule}[\status{D}]\label{sr:distinctness-by-endpoint}
Two branches are distinct only to the extent that certified
endpoint-visible structure distinguishes them.
\end{standingrule}

%% ---------------------------------------------------------------
\section{Branch Candidates and Descent Structure}
%% ---------------------------------------------------------------

\begin{definition}[\status{D}]\label{def:target-category}
A \emph{licensed target category} $\mathcal{C}$ is a mathematical
category to which a faithful realization functor from $\mathrm{POT}$
may be directed, provided the functor is licensed by an explicit
bridge theorem.  Examples include: finite algebraic sector categories,
operator/semigroup categories, and continuum realization categories
(the last requiring the interface legitimacy of Book~VI).
\end{definition}

\begin{definition}[\status{D}]\label{def:realization-functor}
A \emph{realization functor} is a functor $F: \mathrm{POT} \to
\mathcal{C}$ to a licensed target category that: preserves the
POT-fibration structure; maps POT-morphisms to morphisms of
$\mathcal{C}$; and is licensed by an explicit bridge theorem.
\end{definition}

\begin{definition}[\status{D}]\label{def:branch-candidate}
A \emph{licensed branch candidate} is a pair $(\mathcal{C}, F)$
consisting of a licensed target category $\mathcal{C}$ and a
proposed realization functor $F: \mathrm{POT} \to \mathcal{C}$,
together with an explicit declaration of the bridge theorem licensing
$F$.
\end{definition}

\begin{definition}[\status{D}]\label{def:endpoint-datum}
The \emph{endpoint datum} of a licensed branch candidate
$(\mathcal{C}, F)$ is the object
\[
  E \;:=\; F(\mathbf{U}) \;\in\; \mathcal{C},
\]
the image of the universal source object under the realization functor.
\end{definition}

\begin{definition}[\status{D}]\label{def:branch-projection}
A \emph{branch projection} is the canonical morphism
\[
  p_F: \mathbf{U} \to F^{-1}(E)
\]
expressing how $\mathbf{U}$ maps to the pre-image of the endpoint
datum under $F$.  At the POT level, it is the inclusion map
$\eta_R: F_R \to \mathbf{U}$ post-composed with the realization functor.
\end{definition}

\begin{definition}[\status{D}]\label{def:branch-visible}
A \emph{branch-visible observable} is any descendant observable
$\varphi$ on $E \in \mathcal{C}$ whose definability and value are
invariant under replacement by source-equivalent data: if $\mathbf{U}
\cong \mathbf{U}'$ as POT objects, then $F(\mathbf{U}) \cong
F(\mathbf{U}')$, and $\varphi$ assigns the same value to both.
\end{definition}

\begin{definition}[\status{D}]\label{def:legitimacy-witness}
A \emph{legitimacy witness} for a licensed branch candidate
$(\mathcal{C}, F)$ is the certified evidence that:
\begin{enumerate}[label=\textup{(\arabic*)}]
  \item the endpoint datum $E = F(\mathbf{U})$ arises by certified
        descent from $\mathbf{U}$ through the declared realization
        functor $F$;
  \item $E$ is branch-visible: its definition depends only on
        $\mathbf{U}$ and the declared functor $F$; and
  \item no undeclared supplementary structure beyond $\mathbf{U}$, $F$,
        and the declared branch-specific hypotheses is needed to define
        $E$.
\end{enumerate}
\end{definition}

\begin{definition}[\status{D}]\label{def:hidden-channel}
A \emph{hidden channel} is any undeclared dependence of an endpoint
claim on structure not visible through the declared branch projection
and certified realization functor.  Specifically, a hidden channel
exists whenever:
\begin{enumerate}[label=\textup{(\arabic*)}]
  \item the endpoint $E$ requires structure beyond what $F(\mathbf{U})$
        carries; or
  \item the legitimacy of $E$ implicitly assumes a result not in the
        declared import ledger.
\end{enumerate}
\end{definition}

\begin{remark}[\status{R}]\label{rmk:hidden-channel-examples}
Hidden channels are the primary failure mode of ill-disciplined
downstream programs.  Common instances include: silently assuming a
continuum structure not licensed by Book~VI, implicitly importing a
prior result not declared in the branch's import ledger, treating a
conditional result as unconditional, and reading downstream physical
interpretation back into source-level structure.  Standing
Rule~\ref{sr:no-hidden-supplementation} prohibits all such imports
regardless of how they are framed.
\end{remark}

%% ---------------------------------------------------------------
\section{Branch Legitimacy}
%% ---------------------------------------------------------------

\begin{proposition}[\status{U}]\label{prop:branch-projection}
\textup{(Branch Projection Theorem.)}  For every licensed branch
candidate $(\mathcal{C}, F)$, the endpoint datum $E = F(\mathbf{U})$
is well-defined in $\mathcal{C}$, and the branch projection is
canonical: it depends only on $\mathbf{U}$ and the declared
realization functor $F$.
\end{proposition}

\begin{proof}
$\mathbf{U}$ exists and is unique up to unique POT-isomorphism
(Book~IV, Theorems~IV.4.1 and IV.5.1).  A realization functor
$F: \mathrm{POT} \to \mathcal{C}$ maps POT-isomorphisms to isomorphisms
in $\mathcal{C}$.  Therefore $E = F(\mathbf{U})$ is well-defined up to
isomorphism in $\mathcal{C}$, and the branch projection is canonical.
\end{proof}

\begin{proposition}[\status{U}]\label{prop:endpoint-visibility}
Under a licensed realization functor $F$, the endpoint datum $E =
F(\mathbf{U})$ is branch-visible: its observational content is
invariant under replacement by source-equivalent data.
\end{proposition}

\begin{proof}
If $\mathbf{U} \cong \mathbf{U}'$ as POT objects (i.e., they are
source-equivalent), then $F(\mathbf{U}) \cong F(\mathbf{U}')$ in
$\mathcal{C}$ since $F$ preserves isomorphisms.  Therefore the
observational content of $E$ is invariant under source equivalence.
Branch-visible observables on $E$ inherit this invariance.
\end{proof}

\begin{proposition}[\status{U}]\label{prop:legitimacy-witness-required}
\textup{(Legitimacy Witness Requirement.)}  A proposed downstream
branch has canonical endpoint force only if it admits a legitimacy
witness.
\end{proposition}

\begin{proof}
By Standing Rule~\ref{sr:source-descent-only} and
Definition~\ref{def:legitimacy-witness}: canonical endpoint force
requires that the endpoint datum depend only on the declared source
descent and declared branch-specific hypotheses.  Without a legitimacy
witness certifying this dependency, the endpoint claim lacks the
anti-smuggling guarantee required for canonical force.
\end{proof}

\begin{proposition}[\status{U}]\label{prop:hidden-channel-exclusion}
\textup{(Hidden Channel Exclusion.)}  Any proposed downstream branch
depending on a hidden channel is non-canonical.  Specifically, if the
endpoint datum $E$ requires undeclared structure beyond $F(\mathbf{U})$
and the declared hypotheses, then no legitimacy witness exists.
\end{proposition}

\begin{proof}
A legitimacy witness certifies (condition~(3) of
Definition~\ref{def:legitimacy-witness}) that no undeclared
supplementary structure is needed.  If a hidden channel exists
(Definition~\ref{def:hidden-channel}), then condition~(3) fails:
there is undeclared structure influencing the endpoint datum.
Therefore the legitimacy witness condition cannot be satisfied, and
the branch claim is non-canonical.
\end{proof}

\begin{corollary}[\status{U}]\label{cor:hidden-channel-non-canonical}
Any descendant relying on hidden-channel supplementation is
non-canonical.
\end{corollary}

\begin{remark}[\status{R}]\label{rmk:hidden-channel-anti-smuggling}
Proposition~\ref{prop:hidden-channel-exclusion} is the descendant-level
anti-smuggling hinge of Book~V.  A branch becomes lawful not when it
sounds plausible, but when its endpoint can be certified as
source-descended without hidden supplementation.  The three conditions
of the legitimacy witness --- certified descent, branch visibility,
and no undeclared supplementation --- together constitute the canonical
test for branchhood.
\end{remark}

%% ---------------------------------------------------------------
\section{Distinctness and Collapse Discipline}
%% ---------------------------------------------------------------

Lawful descent does not by itself decide when two branches are
genuinely distinct.  Book~V must distinguish route-difference from
endpoint-difference.

\begin{definition}[\status{D}]\label{def:path-equiv}
Two branch projections $p_F$ and $p_{F'}$ are \emph{path-equivalent}
if they arise from the same declared realization functor $F = F'$ via
the same certified descent route from $\mathbf{U}$.
\end{definition}

\begin{definition}[\status{D}]\label{def:endpoint-equiv}
Two branch projections $p_F$ and $p_{F'}$ are \emph{endpoint-equivalent}
if their endpoint data agree up to the native equivalence of their
target categories: $F(\mathbf{U}) \cong F'(\mathbf{U})$ in the
relevant licensed category.
\end{definition}

\begin{definition}[\status{D}]\label{def:collapse-case}
A \emph{collapse case} occurs when two proposed descendants are both
path-equivalent and endpoint-equivalent.
\end{definition}

\begin{definition}[\status{D}]\label{def:confluent-case}
A \emph{confluent distinctness case} occurs when two descendants are
path-distinct but endpoint-equivalent: they arrive at the same
certified endpoint structure via genuinely different certified routes.
\end{definition}

\begin{definition}[\status{D}]\label{def:fully-distinct}
A \emph{fully distinct case} occurs when two descendants differ both
by path and by endpoint datum.
\end{definition}

\begin{theorem}[\status{U}]\label{thm:collapse}
\textup{(Collapse Criterion.)}  If two proposed descendants of
$\mathbf{U}$ are path-equivalent and endpoint-equivalent, then they do
not constitute distinct lawful branches.
\end{theorem}

\begin{proof}
If the descendants share the same certified descent route and terminate
in equivalent endpoint data, they differ neither by source path nor by
target content.  By Standing Rule~\ref{sr:distinctness-by-endpoint},
two branches are distinct only to the extent that certified
endpoint-visible structure distinguishes them.  Since no such
distinction is present, the two proposed descendants are canonically
the same branch.
\end{proof}

\begin{corollary}[\status{U}]\label{cor:path-mult-no-branch}
Path multiplicity without endpoint distinction does not by itself
generate new canonical branches.
\end{corollary}

\begin{corollary}[\status{U}]\label{cor:redescription-no-branch}
Endpoint redescription without source-distinct descent does not
generate new canonical branches.
\end{corollary}

\begin{remark}[\status{R}]\label{rmk:collapse-inflection}
Corollaries~\ref{cor:path-mult-no-branch} and
\ref{cor:redescription-no-branch} block the two main failure modes of
branch inflation: constructing a ``new branch'' by simply relabeling
the same descent path, or repackaging the same endpoint datum with
different notation.  Both are excluded by the Collapse Criterion.
\end{remark}

But collapse is not the only possibility.  Distinct lawful paths may
converge at the same endpoint without becoming identical as routes.

\begin{definition}[\status{D}]\label{def:diamond-confluence}
A \emph{diamond confluence configuration} is a pair of lawful descent
paths from $\mathbf{U}$ into the same licensed target category whose
endpoint data agree up to the native equivalence of that category.
\end{definition}

\begin{theorem}[\status{U}]\label{thm:diamond}
\textup{(Diamond Confluence Principle.)}  If two lawful descent paths
from $\mathbf{U}$ into a declared target category are
endpoint-equivalent, then they are canonically confluent at the
endpoint, even when they remain path-distinct.
\end{theorem}

\begin{proof}
Endpoint-equivalence identifies the terminal descendant structure up to
the target category's native equivalence.  Path-distinctness records
that the certified descent routes are different.  Both facts can hold
simultaneously: the endpoint-equivalence does not force the routes
themselves to coincide.  Therefore confluence at the endpoint does not
require collapse of the routes.  Two route-distinct but
endpoint-equivalent descendants remain canonically meaningful as
distinct paths to the same endpoint.
\end{proof}

\begin{corollary}[\status{U}]\label{cor:confluence-no-collapse}
Confluence at the endpoint does not erase lawful path distinction.
\end{corollary}

\begin{corollary}[\status{U}]\label{cor:route-distinct-meaningful}
Distinct descent paths may remain mathematically meaningful even when
they terminate in the same certified endpoint structure.
\end{corollary}

\begin{remark}[\status{R}]\label{rmk:diamond-permits-multiplicity}
Taken together, the Collapse Criterion and the Diamond Confluence
Principle say: Book~V permits lawful multiplicity (distinct routes to
a common endpoint) while forbidding arbitrary branch inflation (calling
the same thing new merely by redescription).  This is the correct
balance.
\end{remark}

%% ---------------------------------------------------------------
\section{The Universal Branch Legitimacy Theorem}
%% ---------------------------------------------------------------

The previous sections establish all five conditions needed for the
governing theorem of this book.

\begin{theorem}[\status{U}]\label{thm:UBLT}
\textup{(Universal Branch Legitimacy Theorem.)}  A proposed downstream
branch $(\mathcal{C}, F, E)$ is lawful if and only if all five of the
following conditions hold:
\begin{enumerate}[label=\textup{(\arabic*)}]
  \item\label{cond:intake} \emph{Licensed intake}: $(\mathcal{C}, F)$
        is a licensed branch candidate with explicitly declared
        realization functor.
  \item\label{cond:descent} \emph{Certified source descent}: the
        endpoint datum $E = F(\mathbf{U})$ arises as the certified
        image of the universal source object $\mathbf{U}$ under the
        declared realization functor.
  \item\label{cond:visibility} \emph{Endpoint visibility}: $E$ is
        branch-visible --- its definition and value are invariant under
        source equivalence.
  \item\label{cond:witness} \emph{Legitimacy witness}: a legitimacy
        witness for $(\mathcal{C}, F)$ exists, certifying that $E$
        depends only on $\mathbf{U}$, the declared functor $F$, and
        the declared branch-specific hypotheses.
  \item\label{cond:no-channel} \emph{No hidden-channel supplementation}:
        the endpoint claim requires no undeclared structure beyond what
        conditions~\ref{cond:intake}--\ref{cond:witness} license.
\end{enumerate}
\end{theorem}

\begin{proof}
\emph{Necessity.}  Each condition is necessary for canonical
branchhood:
\begin{itemize}
  \item Condition~\ref{cond:intake} is necessary by
        Definition~\ref{def:branch-candidate}: a branch must have a
        declared and licensed descent mechanism.
  \item Condition~\ref{cond:descent} is necessary by
        Book~IV, Corollary~IV.5.3: every lawful downstream realization
        factors through $\mathbf{U}$.
  \item Condition~\ref{cond:visibility} is necessary by
        Standing Rule~\ref{sr:endpoint-visibility}: a branch endpoint
        must be branch-visible.
  \item Condition~\ref{cond:witness} is necessary by
        Proposition~\ref{prop:legitimacy-witness-required}: canonical
        force requires a legitimacy witness.
  \item Condition~\ref{cond:no-channel} is necessary by
        Proposition~\ref{prop:hidden-channel-exclusion}: any hidden
        channel invalidates the legitimacy witness.
\end{itemize}

\emph{Sufficiency.}  Conversely, conditions~\ref{cond:intake}--\ref{cond:no-channel}
are jointly sufficient:
\begin{itemize}
  \item Condition~\ref{cond:intake} provides the declared descent
        structure.
  \item Condition~\ref{cond:descent} ensures the endpoint is sourced
        from $\mathbf{U}$.
  \item Condition~\ref{cond:visibility} ensures the endpoint is
        canonically meaningful.
  \item Condition~\ref{cond:witness} certifies the dependency structure
        is free of smuggling.
  \item Condition~\ref{cond:no-channel} confirms no undeclared
        supplementation is present.
\end{itemize}
Once all five conditions are met, the descendant is fully
source-certified, endpoint-visible, and free of hidden support.  No
further conditions are needed for the branch to count as a lawful
descendant.
\end{proof}

\begin{corollary}[\status{U}]\label{cor:mandatory-declarations}
Every derived branch book must begin by declaring:
\begin{enumerate}[label=\textup{(\arabic*)}]
  \item its licensed target category $\mathcal{C}$,
  \item its realization functor $F: \mathrm{POT} \to \mathcal{C}$
        and the bridge theorem licensing it,
  \item its endpoint datum $E = F(\mathbf{U})$,
  \item the native equivalence of $\mathcal{C}$ determining when two
        endpoint data coincide, and
  \item its legitimacy witness.
\end{enumerate}
A branch book that fails to supply these five items is not
canonically admissible under Book~V.
\end{corollary}

\begin{proof}
The five items correspond precisely to the five conditions of
Theorem~\ref{thm:UBLT}.  Admissibility requires all five; omitting any
one leaves the branch without canonical force.
\end{proof}

\begin{corollary}[\status{U}]\label{cor:open-branch-still-canonical}
A branch may be mathematically serious and canonically admissible even
when its closure remains open, provided its target, descent law, and
failure frontier are canonically declared.
\end{corollary}

\begin{proof}
The Universal Branch Legitimacy Theorem certifies that a branch is
\emph{lawful}, not that it is \emph{closed}.  Lawfulness is a
structural condition about source descent and hidden-channel
exclusion.  Closure is a further mathematical achievement requiring
additional branch-specific theorems.  A lawful but unclosed branch
with an explicitly declared failure frontier satisfies all conditions
of Theorem~\ref{thm:UBLT} and retains full canonical standing.
\end{proof}

\begin{scholium}[\status{R}]\label{sch:UBLT-scope}
Book~V defines lawful descendant structure.  It does not yet determine
the exact force of descendant claims, nor does it license continuum
language within them.  The force of a branch claim is governed by
Book~VII.  The licensing of continuum and differential language within
a branch is governed by Book~VI.  Both of those books are downstream
of Book~V: they presuppose that a branch is lawful before governing
how its claims should be stated.
\end{scholium}

%% ---------------------------------------------------------------
%% ---------------------------------------------------------------
\section{Screened Sector as Canonical Endpoint Datum}
\label{sec:screening-inevitability}
%% ---------------------------------------------------------------

This section records the Screening Inevitability Chain
(Holding BIN CROSS, v7.35 §11791) and the canonical
identification of the screened sector as a branch-visible
endpoint datum satisfying Book~V discipline.

\begin{proposition}[\status{C}]\label{prop:tau-channel-ceiling}
\textup{($\tau$-Channel Ceiling Bound.)}
\textup{[Conditional on $K_0=7$, LS-2, Phase-A.]}
The number of independent boundary-channel parameters consistent
with Phase-A stability is bounded above by $\tau(7) = 4$
(three $d_{S_v}$-block directions plus one gauge-invariant
direction).  Any observable family with more than $\tau(7)$
independent boundary-channel parameters violates Phase-A stability.
\end{proposition}

\begin{proof}
By the self-consistency fixed-point theorem (dream\_paper\_v38
§6.3, Theorem~6.5): the number of independent boundary parameters
consistent with Phase-A is $\tau(p) \leq 5$ for all primes (the
Avoidance-Cycle Exclusion theorem caps $\tau$), and for $p=7$
specifically $\tau(7) = 4$ (verified in §6.2).  Any observable
family exceeding this bound cannot arise from Phase-A data. \qed
\end{proof}

\begin{lemma}[\status{C}]\label{lem:no-hidden-resonances}
\textup{(PASS + LS-2 Forbid Hidden Decoupled Near-Resonances.)}
\textup{[Conditional on PASS (Thm.~\ref{thm:PASS-derived},
GR branch), LS-2 (Book~II).]}
Under PASS and LS-2, no decoupled near-resonant boundary mode
can persist in the observable family without being captured by
the screened sector $W_A$: any mode that commutes with the
full $G$-action and survives LS-2 stabilization is already
in $W_A$.
\end{lemma}

\begin{proof}
By PASS (Theorem~\ref{thm:PASS-derived}): every backbone
observable in $W_A$ commutes with $G$.  Conversely, any
$G$-commuting observable not already in $W_A$ would define an
additional $G$-invariant sector not accounted for in the
$d_{S_v}$-block decomposition.  By the $\tau$-channel ceiling
(Proposition~\ref{prop:tau-channel-ceiling}), no such additional
sector exists at $K_0=7$.  Under LS-2, stabilized local
observables are determined by collar type; any near-resonant
mode that is LS-2 stable and $G$-commuting is already captured
by the $K_0 = 7$ collar type structure and hence in $W_A$. \qed
\end{proof}

\begin{theorem}[\status{C}]\label{thm:screened-sector-endpoint}
\textup{(Screened Sector as Canonical Branch-Visible Endpoint Datum.)}
\textup{[Conditional on PASS, LS-2, Phase-A, $K_0=7$.]}
The screened sector $W_A = \ker(T_\partial - \lambda_{\max}I)^\perp$
is a canonical branch-visible endpoint datum in the sense of
Definition~\ref{def:endpoint-datum}:
\begin{enumerate}[label=\textup{(\arabic*)}]
  \item \emph{Source-descended}: $W_A$ is determined by
    $T_\partial$ and the canonical screening rule, both
    source-descended (Lemma~T3.1a, YM Branch).
  \item \emph{Branch-visible}: $W_A$ is the declared endpoint
    of the YM, NS, and SM branches --- the screened observable
    algebra acts on $W_A$ and no hidden channel contributes
    (Lemma~\ref{lem:no-hidden-resonances}).
  \item \emph{Sector-unique}: $W_A$ is the unique complement
    of the dominant eigenspace of $T_\partial$; no alternative
    screened sector exists at the same $\tau$-level
    (Proposition~\ref{prop:tau-channel-ceiling}).
  \item \emph{Book V compliant}: $W_A$ arises as the certified
    image $F(\mathbf{U})|_{W_A}$ of the universal completion
    $\mathbf{U}$ through the canonical screening functor,
    satisfying the five UBLT conditions of
    Definition~\ref{def:endpoint-datum}.
\end{enumerate}
\end{theorem}

\begin{proof}
(1) Lemma~T3.1a (YM Branch: canonical screened-sector definability):
$W_A$ is determined by $(\Sigma_\partial, T_\partial)$ alone,
both source-descended.
(2) By Lemma~\ref{lem:no-hidden-resonances}: all $G$-commuting
LS-2-stable observables are in $W_A$; no hidden channel.
(3) By Proposition~\ref{prop:tau-channel-ceiling}: $\tau(7)=4$
caps the number of independent sectors; $W_A$ is the unique
one matching the declared Phase-A architecture.
(4) The screening functor $\Pi_A : \mathbf{U} \to W_A$
(the canonical orthogonal projection) is a morphism of
the POT category (Book~IV); $W_A = F(\mathbf{U})|_{W_A}$
arises as the screened sub-object.  The five UBLT conditions
follow from (1)--(3) and the lawfulness of $\Pi_A$
(Lemma~T3.1c, YM Branch: bracket-compatible). \qed
\end{proof}

\begin{remark}[\status{R}]
Theorem~\ref{thm:screened-sector-endpoint} closes the Holding
entry ``Screening Inevitability Chain'' at the level of a
conditional canonical theorem.  The one remaining open hinge
--- Uniform Anti-Chasing (the analytic gap in the screening
inevitability argument for non-abelian Gribov copies) ---
is not needed for the endpoint-datum identification theorem;
it would be required for a stronger theorem asserting that
\emph{all} potential hidden channels are provably absent.
\end{remark}

\section{Boundary of Book V}
%% ---------------------------------------------------------------

\textbf{What Book~V proves.}
\begin{enumerate}
  \item Lawful descent from $\mathbf{U}$ is governed by a five-condition
        biconditional: the Universal Branch Legitimacy Theorem.
  \item Every lawful branch must have a licensed target category,
        a declared realization functor, a certified endpoint datum, a
        legitimacy witness, and no hidden-channel supplementation.
  \item Path multiplicity without endpoint distinction does not
        generate new branches (Collapse Criterion).
  \item Distinct descent routes may converge at the same endpoint
        without collapsing (Diamond Confluence Principle).
  \item A branch is canonically admissible even when unclosed,
        provided its failure frontier is explicitly declared.
\end{enumerate}

\textbf{What Book~V does not prove, and may not be assumed here.}
\begin{enumerate}
  \item \emph{Branch closures.}  Book~V certifies branchhood, not
        closure.  Yang--Mills closure, Navier--Stokes regularity, and
        gravity closure all belong in derived branch books.
  \item \emph{Continuum interface.}  Book~V does not license the use
        of differential, integral, geometric, or smooth language in
        branch books.  That licensing belongs to Book~VI.
  \item \emph{Force of branch claims.}  Book~V does not determine
        whether a branch-level claim is unconditional, conditional, or
        bridge-dependent.  That governance belongs to Book~VII.
  \item \emph{Physical endpoint identification.}  The identification
        of any specific endpoint datum with a physical theory (gauge
        field, fluid, spacetime geometry) belongs in derived branch
        books and requires specific bridge theorems not proved here.
\end{enumerate}

%% ---------------------------------------------------------------
\section{Handoff to Book VI and Book VII}
%% ---------------------------------------------------------------

Book~V branches the canonical program in two directions.

\emph{Toward Book~VI:} When a lawful branch seeks to use continuum,
differential, integral, or geometric language as an encoding of its
endpoint structure, it must invoke Book~VI's interface legitimacy
machinery.  No smooth language is licensed in a branch book without
first establishing the relevant continuum interface regime through
Book~VI.

\emph{Toward Book~VII:} When a lawful branch seeks to state the exact
force, scope, import discipline, and failure frontier of its claims,
it must invoke Book~VII's canonical closure schema.  The governance of
how branch claims may be promoted, cited, and combined belongs
constitutionally to Book~VII.

The handoff from Book~V is therefore:
\begin{quote}
  $\mathbf{U}$
  $\;\longrightarrow\;$
  licensed branch candidate
  $\;\longrightarrow\;$
  lawful branch projection
  $\;\longrightarrow\;$
  $\begin{cases}
    \text{Book~VI: interface legitimacy,}\\
    \text{Book~VII: closure discipline.}
  \end{cases}$
\end{quote}
Book~V supplies the first three stages.

%% ---------------------------------------------------------------
\section{Dependency Ledger}
%% ---------------------------------------------------------------

\renewcommand{\arraystretch}{1.4}
\noindent
\begin{tabular}{@{}p{4.2cm}p{0.8cm}p{3.4cm}p{4.6cm}@{}}
\hline
\textbf{Theorem} & \textbf{St.} & \textbf{Depends on} & \textbf{Exports to} \\
\hline
\ref{prop:branch-projection}\ Branch Projection Thm.
  & [U] & Bk.~IV Thms.~IV.4.1, IV.5.1
  & \ref{prop:endpoint-visibility}, \ref{thm:UBLT} \\

\ref{prop:endpoint-visibility}\ Endpoint Visibility
  & [U] & \ref{prop:branch-projection}
  & \ref{thm:UBLT} \\

\ref{prop:legitimacy-witness-required}\ Legitimacy Witness Req.
  & [U] & SR~\ref{sr:source-descent-only},
    Def.~\ref{def:legitimacy-witness}
  & \ref{prop:hidden-channel-exclusion}, \ref{thm:UBLT} \\

\ref{prop:hidden-channel-exclusion}\ Hidden Channel Exclusion
  & [U] & Def.~\ref{def:hidden-channel},
    \ref{prop:legitimacy-witness-required}
  & \ref{cor:hidden-channel-non-canonical}, \ref{thm:UBLT} \\

\ref{thm:collapse}\ Collapse Criterion
  & [U] & Defs.~\ref{def:path-equiv}--\ref{def:collapse-case},
    SR~\ref{sr:distinctness-by-endpoint}
  & \ref{cor:path-mult-no-branch}, \ref{cor:redescription-no-branch} \\

\ref{thm:diamond}\ Diamond Confluence Principle
  & [U] & Defs.~\ref{def:confluent-case}, \ref{def:diamond-confluence}
  & \ref{cor:confluence-no-collapse}, \ref{cor:route-distinct-meaningful} \\

\ref{thm:UBLT}\ Universal Branch Legitimacy Thm.
  & [U] & \ref{prop:branch-projection}--\ref{prop:hidden-channel-exclusion};
    Bks.~I--IV
  & Books~VI--VII; all branch books \\

\ref{cor:mandatory-declarations}\ Mandatory Declarations
  & [U] & \ref{thm:UBLT}
  & Constitutional rule for all branch books \\

\ref{cor:open-branch-still-canonical}\ Open branch canonical
  & [U] & \ref{thm:UBLT}
  & NS, YM, GR, SCC, A* branch books \\
\hline
\end{tabular}
\renewcommand{\arraystretch}{1}

\medskip

\begin{scholium}[\status{R}]\label{sch:V-discipline}
The Universal Branch Legitimacy Theorem is the constitutional entry
point for every derived branch book.  No branch may begin its work
without satisfying the five conditions of Theorem~\ref{thm:UBLT},
declared explicitly at the opening of the branch manuscript.  This
requirement is not bureaucratic: each condition rules out a specific
form of mathematical dishonesty.  Condition~\ref{cond:intake} rules
out unlicensed realization.  Condition~\ref{cond:descent} rules out
branches that do not actually descend from $\mathbf{U}$.
Condition~\ref{cond:visibility} rules out unobservable endpoints.
Condition~\ref{cond:witness} rules out undeclared dependencies.
Condition~\ref{cond:no-channel} rules out hidden imports.  A branch
that satisfies all five conditions has genuinely earned its canonical
status.
\end{scholium}


%% ---------------------------------------------------------------
\section{Coercive-Transport Source Theorem and ICDC Schema}
\label{sec:br-source}
%% ---------------------------------------------------------------

\begin{theorem}[\status{C}]\label{thm:BR-source}
\textup{(Unified Coercive-Transport Source Theorem, BR.)}
\textup{[Conditional on $K_0 = 7$, the universal object
$\mathbf{U} \in \mathrm{POT}$ (Book~IV, Thm.~\ref{thm:universal-completion}),
and the controlled continuum realization
(Book~VI, Thm.~\ref{thm:continuum-bridge-schema}).]}
There exists a single upstream coercive-transport structure on
$\mathbf{U}$ with the following burden-routed content:
\begin{enumerate}[label=\textup{(\roman*)}]
  \item \emph{Source-level existence.}  A coercive form
    $\mathcal{B}(f,g) = \int_{\mathbf{U}} \langle\nabla_{\mathbf{U}}
    f, \nabla_{\mathbf{U}} g\rangle\,d\mu_{\mathrm{nl}}$ is
    well-defined on $\mathcal{H}_\infty = L^2(\mathbf{U}, \mu_{\mathrm{nl}})$,
    with screened-sector lower bound
    $\mathcal{B}(f,f) \geq c\|f\|^2$ for some $c > 0$.
    This is the intrinsic content of the BR theorem.
  \item \emph{Yang--Mills specialization.}  Under the realization
    and quantization inputs of the YM Branch, the source coercivity
    pushes forward to the Yang--Mills branch Hamiltonian estimate.
  \item \emph{Navier--Stokes specialization.}  Under the realization
    and transport inputs of the NS Branch, the same source pushes
    forward to the NS branch floor/transport estimate.
  \item \emph{Structural Einstein specialization.}  Under the
    geometric realization inputs of the GR Branch (KPO-3), the same
    source pushes forward to the EFE coercive structure.
\end{enumerate}
Accordingly, the YM, NS, and GR branch theorems are to be read as
certified specializations of one upstream source theorem, not as
independent foundations.
\end{theorem}

\begin{proof}
The unconditional existence of the coercive form follows from
Thm.~\ref{thm:universal-completion} (Book~IV), which provides both
the transfer structure on $\mathbf{U}$ and the canonical transport
measure $\mu_{\mathrm{nl}}$.  Branch specialization under $K_0 = 7$
follows from the controlled continuum realization
(Book~VI, Thm.~\ref{thm:continuum-bridge-schema}), which licenses
the push-forward of the coercive structure to the realized operators
in each branch.  No branch conclusion belongs to the unconditional BR
theorem beyond this burden-routed transfer principle.
(Attributed: dream\_paper\_v38.pdf Thm.~9.19.) \qed
\end{proof}

\begin{theorem}[\status{C}]\label{thm:ICDC}
\textup{(Interface-Coherent Directed Compression, ICDC Schema.)}
\textup{[Conditional on $K_0 = 7$ and one-step inheritance.]}
Let $U_{n_0} \subseteq U_{n_0+1} \subseteq \cdots$ be an admissible
enlargement ladder from $\mathbf{U}$, and let $(B_n)_{n \geq n_0}$ be
a branch family functorially extracted along this ladder.  Assume:
\begin{enumerate}[label=\textup{(\arabic*)}]
  \item \emph{One-step inheritance}: branch data transport lawfully
    from $B_n$ to $B_{n+1}$;
  \item \emph{Interface localization}: any failure of exact
    inheritance is confined to an explicit interface residue
    $\partial_n$;
  \item \emph{Two-step transitivity}: corrections for $n \to n+1
    \to n+2$ compose into a single accumulated residue $\partial_{n,n+2}$;
  \item \emph{Source-image discipline}: every load-bearing datum of
    each $B_n$ remains in the image of $\mathbf{U}$ with no new
    branch-local primitive introduced;
  \item \emph{Admissible compression}: the accumulated interface
    residues are summably admissible or uniformly subordinate to a
    branch-specific coercive margin.
\end{enumerate}
Then the family admits a controlled scale-limit compression
$B_\infty := \varinjlim_\partial B_n$ in which the load-bearing
branch data survive coherently across scale modulo explicit interface
bookkeeping.
\end{theorem}

\begin{proof}
The one-step inheritance maps define the direct-system spine.
Interface localization isolates all failure of exact transport in
named residues.  Two-step transitivity ensures residues assemble
coherently.  Source-image discipline blocks hidden downstream
smuggling.  Admissible compression yields a well-defined controlled
limit object.
(Attributed: dream\_paper\_v38.pdf Thm.~9.12.) \qed
\end{proof}

\begin{remark}[\status{R}]\label{rem:icdc-branches}
\textbf{Branch-specific ICDC readings.}
The ICDC schema specializes to three terminal branches:
\begin{itemize}
  \item \emph{Yang--Mills} ($\Theta_n = 1$, gap-subcritical): gap
    preservation holds whenever $E_n + B_n < g_\infty^{(n)}$;
  \item \emph{Navier--Stokes} ($\Theta_n = \theta < 1$): obstruction
    contraction $C_{n+1}^{NS} \leq \theta C_n^{NS} + E_n + B_n$;
  \item \emph{Gravity} ($\Theta_n = 1$, curvature-subcritical):
    domain extension holds when new interface contributes curvature
    defect strictly subordinate to the already-controlled gravity margin.
\end{itemize}
The common upstream law is the interface-relative stability
functional $\mathcal{C}_n(U_n) = B_n^{\mathrm{bulk}}(U_n) -
\mathcal{I}_n(U_n)$; what changes branch-to-branch is only the
terminal observable and closure criterion.
\end{remark}

\end{document}
